% 392_Darstellungen_Trennung_De_ch.tex
% Johann Pascher, 11. Okt. 2026
% Wo die Darstellungen der Gravitation sich trennen

\providecommand{\BM}{\textbf{[B]}}
\providecommand{\KM}{\textbf{[K]}}
\providecommand{\SM}{\textbf{[S]}}
\providecommand{\XM}{\textbf{[X]}}
\providecommand{\QM}{\textbf{[Q]}}

\chapter*{Wo die Darstellungen der Gravitation sich trennen}
\addcontentsline{toc}{chapter}{Wo die Darstellungen der Gravitation sich trennen}

\begin{abstract}
Dok.~389 führt Raumkrümmung, Massen- bzw.\ Taktänderung und fraktale Wegverlängerung als drei
Darstellungen der Gravitation, die aus der Dualität $\tilde T\cdot m=1$ heraus äquivalent sein müssen.
Dieses Dokument prüft, wie weit die Gleichwertigkeit reicht. Die zeitliche Hälfte
$g_{00}=(m/m_0)^2$ ist in jeder zugelassenen Darstellung enthalten und trägt allein Fall, Rotverschiebung
und Uhrengang. Reine Massenvariation, bei der Takte und Längen gemeinsam mit $1/m$ laufen, ist eine
konforme Skalierung des flachen Raums und hat verschwindenden Weyl-Tensor. Das flache
Friedmann-Universum ist ebenfalls konform flach; dort ist die Massenlesart exakt und besteht
Lichtkurven, Hintergrundtemperatur und Tolman-Test, während reine Wegverlängerung am Photon an allen
dreien scheitert. Das Feld eines Sterns trägt dagegen Weyl-Krümmung, deren Komponente
$C_{trtr}=-2GM/r^3$ der Gezeitenterm ist. Dort braucht die flache Darstellung die räumliche Hälfte mit
umgekehrtem Vorzeichen der Maßstäbe, für rotierende Massen ein mitgeführtes Vektorfeld und für
Gravitationswellen einen Tensoranteil. Die Grenze der Lichtgeschwindigkeit folgt aus dem Torusspektrum: Masselose Moden laufen mit $c$, massive
sind immer langsamer; eine konforme Skalierung ändert $c$, $\varepsilon_0$ und $\mu_0$ lokal nicht.
Die Massenlesart ist damit die vollständige Darstellung des
konform flachen Anteils; der Weyl-Anteil ist das, was zusätzlich zu liefern ist. Der
Alcock--Paczyński-Test trennt Expansion und Massenlesart nicht, schließt aber die reine
Exponentialform $1+z=e^{\xi x}$ des Entfernungsgesetzes mit den DESI-Daten aus.
\end{abstract}

% ============================================================
\section*{Statusmarkierungen}
% ============================================================
\begin{center}
\begin{tabular}{@{}lp{11cm}@{}}
\textbf{[B]} & Algebraisch bewiesen (symbolische Rechnung im Prüfskript). \\[2pt]
\textbf{[K]} & Numerisch bestätigt (Prüfskript, Vergleich mit Messwerten). \\[2pt]
\textbf{[S]} & Setzung, offen oder spekulativ. \\[2pt]
\textbf{[X]} & Durch Messung ausgeschlossen. \\[2pt]
\textbf{[Q]} & Aus der Literatur übernommen, mit Quelle. \\
\end{tabular}
\end{center}
Prüfskripte: \texttt{2/python/Dok392\_Skripte/} (acht Skripte, 81/81). Messwerte und Quellen stehen
in den Skriptköpfen und im Literaturverzeichnis.

% ============================================================
\section*{Formelzeichen}
% ============================================================
\begin{center}
\small
\begin{tabular}{@{}p{3.3cm}p{11.2cm}@{}}
\toprule
$\xi=4/30000$ & einziger Parameter der FFGFT (Geometrie T$^4$/Z$_3$) \\
$\tilde T$, $m$, $m_0$ & Zeitfeld, ortsabhängige Masse, Bezugsmasse; Dualität $\tilde T\cdot m=1$ \\
$c$ & Lichtgeschwindigkeit, in der Grundform $c=1$ \\
$G$, $M$, $r$, $\vec x$ & Gravitationskonstante, Masse des Zentralkörpers, Abstand, Ortsvektor \\
$\Phi$, $U$ & Newtonsches Potential; $U=GM/(rc^2)=-\Phi/c^2$, dimensionslos \\
$g_{\mu\nu}$ & Metrik; $g_{00}$ zeitliche Komponente (Takt), $g_{0i}$ gemischte (Mitführung),
$g_{ij}$ räumliche Komponenten ($i,j=1,2,3$) \\
$\eta_{\mu\nu}$, $\Omega$ & flache Minkowski-Metrik; konformer Faktor, bei reiner Massenvariation
$\Omega=m/m_0$ \\
$\gamma$, $\beta$ & PPN-Parameter: Raumkrümmung je Masse; Nichtlinearität von $g_{00}$ \\
$\eta=4\beta-\gamma-3$ & Nordtvedt-Parameter (nicht zu verwechseln mit $\eta_{\mu\nu}$ und der
konformen Zeit $\eta$) \\
$m_G/m_I$ & Verhältnis von schwerer zu träger Masse (Mondlaser) \\
$C_{trtr}$ & Komponente des Weyl-Tensors, der Gezeitenanteil der Krümmung \\
$n$ & Brechungsindex der Wegverlängerung, $n=1+2U$ \\
$\vec J$, $I$, $R$ & Drehimpuls, Trägheitsmoment und Radius des rotierenden Körpers \\
$\vec u$ & Geschwindigkeitsfeld, mit dem die Wegverlängerung mitgeführt wird \\
$\vec w$ & magnetoelektrische Kopplung in der Medium-Form, $\vec w=-g_{0i}/g_{00}$ \\
$\dot P$ & Änderung der Umlaufperiode eines Doppelsterns \\
$z$, $a$, $a_0$ & Rotverschiebung; Skalenfaktor bzw.\ in der Massenlesart $m/m_0$; heutiger Wert \\
$\eta$ (kosmisch), $H$, $H_0$ & konforme Zeit; Hubble-Rate $H=a'/a^2$; ihr heutiger Wert \\
$b$, $\beta$ (bei $T(z)$), $k$ & Exponenten der Lichtkurvenstreckung $(1+z)^b$, der
Hintergrundtemperatur $T\propto(1+z)^{1-\beta}$ (nicht der PPN-Parameter) und der
Flächenhelligkeit $(1+z)^{-k}$ \\
W, T, M & kosmische Lesarten: nur Photon; zusätzlich Uhren; zusätzlich Uhren und Maßstäbe \\
$d_L$, $d_A$ & Leuchtkraft- und Winkeldurchmesser-Entfernung \\
$D_M$, $D_H$, $r_d$ & quere mitbewegte Entfernung, Hubble-Entfernung $D_H=c/H$, Länge des
Schallhorizonts (BAO-Maßstab) \\
$x$ & Wegstrecke im Gesetz $1+z=e^{\xi x}$ (Dok.~280, 308) \\
$E$, $p$, $v$ & Energie, Impuls, Gruppengeschwindigkeit $v=dE/dp$ einer Mode \\
$x_4$, $R_4$, $k_4$, $\lambda_4$ & Koordinate und Radius der kompakten vierten Richtung, Wicklungszahl
($m=k_4/R_4$), Wellenlänge mit $\lambda_4\cdot m=2\pi$ \\
$\varepsilon$, $\mu$, $\varepsilon_0$, $\mu_0$ & effektive Permittivität und Permeabilität im Feld;
elektrische und magnetische Feldkonstante \\
$A$, $B$ & Metrikfunktionen im statischen Feld, $g_{00}=-A$, $g_{ij}=B\,\delta_{ij}$ \\
$\chi^2$ & Summe der quadrierten Abweichungen, gemessen in Standardabweichungen \\
\bottomrule
\end{tabular}
\end{center}

% ============================================================
\section{Die Frage}
% ============================================================

Dok.~389 liest Gravitation in drei Darstellungen: als Raumkrümmung wie in der Allgemeinen
Relativitätstheorie, als Massen- bzw.\ Taktänderung in flachem Raum (Dok.~078) und als fraktale
Wegverlängerung (Dok.~026, 041, 182, 308). Dort ist bereits gezeigt, dass die Massenänderung allein nur
die zeitliche Hälfte trägt und erst mit der Wegverlängerung die Lichtablenkung der ART ergibt. Offen
blieb, ob die drei Darstellungen über das statische Schwachfeld hinaus gleichwertig sind, und wenn
nicht, wo genau sie sich trennen. Dieses Dokument rechnet beide Seiten nach: was eine Darstellung
ausschließt und was für die Massenlesart spricht.

% ============================================================
\section{Die zeitliche Hälfte als gemeinsamer Kern}
% ============================================================

In der Massendarstellung ist $g_{00}=(m/m_0)^2$ mit $m=m_0\,e^{\Phi/c^2}$ (Dok.~389). Der Fall folgt
exakt aus dem Massengradienten, $\ddot{\vec x}=-c^2\nabla\ln m=-\nabla\Phi$ \BM, und weil
$\nabla\ln m$ für jede Mode derselbe ist, ist das schwache Äquivalenzprinzip eingebaut. Dieselbe
Funktion bestimmt die Rotverschiebung und den Uhrengang. Die Messungen dieser Hälfte sind:
\begin{center}
\begin{tabular}{@{}lll@{}}
\toprule
Messung & Erwartung aus $g_{00}$ & gemessen \\
\midrule
Pound--Rebka, 22{,}5\,m & $2{,}45\times10^{-15}$ & $(2{,}57\pm0{,}26)\times10^{-15}$ \QM \\
optische Uhren, 33\,cm & $3{,}6\times10^{-17}$ & $(4{,}1\pm1{,}6)\times10^{-17}$ \QM \\
GPS, Taktanteil je Tag & $+45{,}7\,\mu$s & im Betrieb korrigiert \\
Gravity Probe A, Galileo 5/6 & Taktanteil & bestätigt auf $7\times10^{-5}$ bzw.\ $2{,}5\times10^{-5}$ \QM \\
\bottomrule
\end{tabular}
\end{center}
\KM. Diese Hälfte ist keine Alternative zur Raumkrümmung, sondern in jeder zugelassenen Darstellung
dieselbe; auch in der ART ist $g_{00}$ der Uhrengang. Die Massenlesart beschreibt sie wörtlich: Im
statischen Feld bleibt die Frequenz des Photons auf dem ganzen Weg konstant, verschieden sind die
Atome an den beiden Orten (Dok.~312). Mit $\tilde T\cdot m=1$ ist das eine einzige Größe.

% ============================================================
\section{Was im Schwachfeld ausgeschlossen ist}
% ============================================================

Im parametrisierten Schwachfeld unterscheiden sich die Darstellungen in $\gamma$ (räumliche Hälfte)
und $\beta$ (zweite Ordnung von $g_{00}$). In der Massendarstellung folgt $\beta$ aus der Form von
$m(\Phi)$: Die exponentielle Form ergibt $g_{00}=1-2U+2U^2$, also $\beta=1$, die lineare
$g_{00}=1-2U+U^2$, also $\beta=\tfrac12$ \BM. Vier unabhängige Messungen trennen:
\begin{center}
\small
\begin{tabular}{@{}lrrcccc@{}}
\toprule
Darstellung & $\gamma$ & $\beta$ & Licht & Merkur & GP-B geodät. & Mondlaser \\
\midrule
Raumkrümmung (ART) & 1 & 1 & $1{,}751''$ & $42{,}98''$ & 6604 & $\eta=0$ \\
Takt (exp.) + Wegverlängerung & 1 & 1 & $1{,}751''$ & $42{,}98''$ & 6604 & $\eta=0$ \\
Takt (linear) + Wegverlängerung & 1 & $\tfrac12$ & $1{,}751''$ & $50{,}1''$ \XM & 6604 & $\eta=-2$ \XM \\
nur Takt & 0 & 1 & $0{,}876''$ \XM & $14{,}3''$ \XM & 2201 \XM & $\eta=1$ \XM \\
konform flach (Dok.~143) & $-1$ & $\tfrac12$ & $0$ \XM & $-7{,}2''$ \XM & $-2201$ \XM & $\eta=0$ \\
\bottomrule
\end{tabular}
\end{center}
\KM. Merkur in Bogensekunden je Jahrhundert, geodätische Präzession in mas je Jahr,
$\eta=4\beta-\gamma-3$. Die Messwerte sind $\gamma-1=(2{,}1\pm2{,}3)\times10^{-5}$ (Cassini), die
geodätische Präzession $6601{,}8\pm18{,}3$\,mas/Jahr (Gravity Probe B) und für den Mondlaser
$(m_G/m_I)_E-(m_G/m_I)_M=(-0{,}8\pm1{,}3)\times10^{-13}$ \QM; mit der Differenz der Eigenenergien von
Erde und Mond folgt daraus $\eta=(1{,}8\pm2{,}9)\times10^{-4}$ ohne Korrektur der Zusammensetzung \KM. Die Massenänderung allein scheitert damit auch
an einem Test ohne Licht, der geodätischen Präzession, mit rund 240 Standardabweichungen. Die
lineare Massenform scheitert unabhängig an Merkur und am Mondlaser; die exponentielle Form, in
Dok.~389 als Setzung geführt, ist damit von zwei Seiten festgelegt.

Mischformen sind zulässig. Trägt die Krümmung den Anteil $w$ und die Wegverlängerung den Anteil $s$
der räumlichen Hälfte, so ist $\gamma=w+s$, und die Bedingung lautet $w+s=1$ \KM. Die Freiheit der
Darstellung liegt also in der Aufteilung der räumlichen Hälfte, nicht in ihrer Größe.

% ============================================================
\section{Konform flach oder Weyl-gekrümmt}
% ============================================================

Laufen Takte und Längen gemeinsam mit $1/m$, so ist die Metrik eine konforme Skalierung des flachen
Raums, $g_{\mu\nu}=\Omega^2\eta_{\mu\nu}$ mit $\Omega=m/m_0$. Für jede solche Metrik verschwindet der
Weyl-Tensor \BM. Reine Massenvariation trägt deshalb genau die konform flache Geometrie, und Licht
läuft in ihr wie im flachen Raum.

\textbf{Das Universum im Großen.} Das flache Friedmann-Universum lässt sich in konformer Zeit als
$a(\eta)^2\eta_{\mu\nu}$ schreiben; auch sein Weyl-Tensor verschwindet \BM. Im statischen Rahmen
laufen die Massen mit $m(\eta)=m_0\,a(\eta)$, und $1+z=m_{\text{heute}}/m_{\text{Emission}}=a_0/a$
(Wetterich~\cite{Wetterich2013}) \QM. Weil die kosmische Geometrie konform flach ist, ist die
Massenlesart dort nicht nur beobachtungsgleich, sondern exakt dieselbe Geometrie in anderen Variablen.

\textbf{Das Feld eines Sterns.} Für das isotrope Schwachfeld mit $\gamma=1$ ergibt die Rechnung in
erster Ordnung
\begin{equation}
C_{trtr}=-\frac{2GM}{r^3}\neq0\quad\BM,
\end{equation}
genau der Gezeitenterm. Im Vakuum außerhalb einer Masse ist die Krümmung reine Weyl-Krümmung. Eine
konforme Skalierung kann sie nicht erzeugen.

\textbf{Das Vorzeichen der Maßstäbe.} Tief im Potential ist $m/m_0=1-U$. Eine Länge, die wie die
Compton- oder Bohr-Länge mit $1/m$ läuft, erscheint vom fernen Bezug aus um $1+U$ länger. Die
räumliche Hälfte mit $\gamma=1$ verlangt dagegen Eigenmaßstäbe, die um $1-U$ kürzer erscheinen
\BM. Die Vorzeichen sind entgegengesetzt; das ist die Gegenläufigkeit, die Dok.~389 für
$\lambda_4\cdot m=2\pi$ festhält. Im Sternfeld folgt die räumliche Hälfte also nicht aus der
Massenskalierung, im Friedmann-Universum dagegen schon.

% ============================================================
\section{Kosmos: welche statische Lesart gleichwertig ist}
% ============================================================

Für die Rotverschiebung ohne Expansion gibt es drei statische Lesarten (Dok.~280, 312, 391): die
Wegverlängerung bzw.\ der Energieverlust nur am Photon (W), dieselbe mit Uhren, die über
$\tilde T\cdot m=1$ mitgebucht werden (T), und die volle Massenlesart, in der Uhren und Maßstäbe
mitgebucht werden (M). Dok.~312 nennt vier Diskriminatoren, deren Vorwärtsnachweis offen war.
Drei davon lassen sich direkt rechnen:
\begin{center}
\begin{tabular}{@{}lccc@{}}
\toprule
& Lichtkurven $(1+z)^b$ & $T(z)\propto(1+z)^{1-\beta}$ & Tolman $(1+z)^{-k}$ \\
\midrule
Messung \QM & $b=1{,}003\pm0{,}011$ & $\beta=-0{,}007\pm0{,}027$ & $k=2{,}59\pm0{,}17$ (R), $3{,}37\pm0{,}13$ (I) \\
W: nur Photon & $b=0$ \XM & $\beta=1$ \XM & $k=1$ \XM \\
T: + Uhren & $b=1$ & $\beta=0$ & $k=2$ \XM \\
M: + Uhren + Maßstäbe & $b=1$ & $\beta=0$ & $k=4$ \\
\bottomrule
\end{tabular}
\end{center}
\KM. Die Lichtkurven schließen die Lesart W mit rund 90 Standardabweichungen aus, die
Hintergrundtemperatur mit rund 37. Beim Tolman-Test liegen die beobachteten Werte über $k=2$, obwohl
die Entwicklung der Sternpopulationen frühe Galaxien heller macht und den wahren Exponenten noch
erhöht; nur $k=4$ ist verträglich. Gleichwertig mit der Expansion ist damit genau die Lesart M. In
ihr werden die Maßstäbe von selbst mitgebucht, weil atomare Längen mit $1/m$ laufen
($\lambda_4\cdot m=2\pi$, Dok.~314). Die Wegverlängerung am Photon allein (Dok.~280) ist als
kosmische Lesart nicht ausreichend; sie gilt als Beschreibung derselben Rotverschiebung innerhalb
der Massenlesart, nicht als eigenständige Alternative.

Der vierte Diskriminator, der Alcock--Paczyński-Test, vergleicht die radiale mit der queren Ausdehnung
eines Objekts. Rechnet man im flachen Raum mit $m(t)=m_0a(t)$, mit Längen und Takten proportional zu
$1/m$ und erhaltener Photonenzahl, ergeben sich Winkeldurchmesser- und Leuchtkraftentfernung sowie das
AP-Verhältnis genau wie im FRW-Modell mit demselben $a$; insbesondere gilt $d_L=(1+z)^2d_A$ \BM. Die
Flächenhelligkeit fällt mit $(d_A/d_L)^2$; der Tolman-Exponent $k=4$ ist also dieselbe Aussage wie
diese Reziprozität, und die Lesarten W und T verletzen sie mit $d_L/d_A=(1+z)^{1/2}$ bzw.\ $(1+z)$. Der
AP-Test trennt deshalb Expansion und Massenlesart nicht; er prüft den Verlauf $H(z)$, also das
Entfernungsgesetz.

Für das Gesetz $1+z=e^{\xi x}$ aus Dok.~280 und 308 lässt sich der Test durchführen, wenn man $x$ als
Koordinatenabstand der Lesart M liest. Dann ist $a(\eta)=e^{H_0(\eta-\eta_0)}$, also $H(z)=H_0(1+z)$,
und das Verhältnis von querer zu radialer Entfernung wird
\begin{equation}
\frac{D_M}{D_H}=(1+z)\ln(1+z).
\end{equation}
Es enthält weder $H_0$ noch die Länge $r_d$ des Schallhorizonts; die Einschränkung P20, die
Dok.~267 für die BAO-Abstände nennt, entfällt damit. DESI DR2~\cite{DESI2025} misst dieses Verhältnis
in sechs Rotverschiebungsbereichen \QM. Bis $z\approx1{,}5$ liegt das Gesetz innerhalb von drei
Standardabweichungen, bei $z=0{,}51$ trifft es den Messwert genau. Beim Lyman-$\alpha$-Wald bei
$z=2{,}33$ ergibt es $4{,}01$ gegen gemessene $4{,}52\pm0{,}10$, also $5{,}3$ Standardabweichungen zu
wenig; über alle sechs Werte ist $\chi^2=37{,}3$, für flaches $\Lambda$CDM mit $\Omega_m=0{,}2975$
dagegen $6{,}0$ \KM. Die reine Exponentialform ist in dieser Lesung damit ausgeschlossen \XM. Die
Entartung, die Dok.~267 für die BAO-Winkel annimmt, setzt eine Beziehung $z(d)$ voraus, die bei hoher
Rotverschiebung vom exponentiellen Gesetz abweicht; eine solche Beziehung ist in der FFGFT nicht
hergeleitet \SM.

Die Gleichwertigkeit allein entscheidet nicht zwischen Expansion und Massenlesart. Für die
Massenlesart spricht die einfachere Setzungsstruktur: Die kosmologische Konstante tritt in ihr als
Lesart-Artefakt auf und wird nicht gesetzt (Dok.~308, 312) \SM.

% ============================================================
\section{Rotierende Massen}
% ============================================================

Die Mitführung durch eine rotierende Masse steckt in $g_{0i}$ (Dok.~312). Ein statisches Paar
skalarer Felder, Masse $m(x)$ und Brechungsindex $n(x)$, hat $g_{0i}=0$: Die Massendichte einer starr
rotierenden Kugel ist stationär wie die einer ruhenden, das Drehmoment $J$ tritt in einem Skalar
nicht auf, und das gravitomagnetische Feld $\nabla\times(\vec J\times\vec x/r^3)$ ist ein Dipolfeld,
das nur ein Vektorfeld tragen kann \BM. Gravity Probe B misst die Mitführung zu
$37{,}2\pm7{,}2$\,mas/Jahr bei einer Vorhersage von rund 39, LAGEOS bestätigt sie auf etwa 10\,\%
\QM; eine fehlende Mitführung ist damit mit 5 bzw.\ etwa 10 Standardabweichungen ausgeschlossen
\KM.

Die flache Darstellung erreicht die Mitführung, wenn die Wegverlängerung wie ein bewegtes Medium
mitgeführt wird. Die Gordon-Metrik eines Mediums mit Index $n$ und Geschwindigkeit $\vec u$ hat
$\tilde g_{0i}=-(1-n^{-2})\,u_i/c\approx-4U\,u_i/c$ für $n=1+2U$; Gleichsetzen mit dem ART-Wert
$g_{0i}=-2G(\vec J\times\vec x)_i/(c^3r^3)$ ergibt
\begin{equation}
\vec u=\frac{\vec J\times\vec x}{2Mr^2}\quad\BM.
\end{equation}
An der Erdoberfläche sind das am Äquator etwa $77$\,m/s, ein Sechstel der Rotationsgeschwindigkeit
($I/(2MR^2)$), bei der Sonne $70$\,m/s oder $3{,}5\,\%$; das Feld fällt mit $1/r$ ab \KM. Dieses
Feld ist ein Vektor und folgt nicht aus $m(x)$. In der Krümmungsdarstellung ist die Mitführung
enthalten, weil die FFGFT die Feldgleichungen übernimmt (Dok.~312); ob die Struktur des Trägers das
Mitführungsfeld der flachen Darstellung selbst liefert, ist offen \SM.

% ============================================================
\section{Gravitationswellen}
% ============================================================

Gravitationswellen stellen drei Anforderungen an jede Darstellung. Erstens laufen sie mit
Lichtgeschwindigkeit: Bei GW170817 kam das Gammasignal $1{,}74$\,s nach den Wellen über rund
$40$\,Mpc an, $|c_{\text{gw}}/c-1|\lesssim10^{-15}$ \KM. Zweitens wirkt die räumliche Hälfte auf Licht
und Wellen gleich: Die Shapiro-Verzögerung im Potential der Milchstraße beträgt rund $240$\,Tage;
wirkte die Wegverlängerung nur auf Licht, käme es etwa $120$\,Tage zu spät an \KM. Drittens
strahlen Doppelsterne Quadrupolstrahlung ab: Die Bahnzerfallsrate des Hulse-Taylor-Pulsars folgt der
Quadrupolformel, $\dot P_{\text{GR}}=-2{,}402\times10^{-12}$, auf $0{,}2\,\%$ \KM. Die Polarisation
bevorzugt reine Tensormoden gegenüber reinen Skalarmoden mit einem Bayes-Faktor über $1000$ \QM.
Ein skalares $m(x)$ kann mit $c$ laufen, liefert aber keine Tensorpolarisation; die flache
Darstellung braucht für Wellen einen Tensoranteil \SM. In der Krümmungsdarstellung ist er mit den
übernommenen Feldgleichungen gegeben (Dok.~312).

% ============================================================
\section{Die Grenze der Lichtgeschwindigkeit und die Feldkonstanten}
% ============================================================

In der Grundform mit $c=1$ ist $c$ der Umrechnungsfaktor zwischen den Raumrichtungen und der
kompakten vierten Richtung $x_4$, keine Geschwindigkeit eines Mediums (Dok.~134). Das Spektrum des
Trägers ist $E^2=p^2+m^2$ mit $m=k_4/R_4$; die Wicklungszahl ist die Massenzahl (Dok.~314). Die
Gruppengeschwindigkeit ist
\begin{equation}
v=\frac{dE}{dp}=\frac{p}{\sqrt{p^2+m^2}}\quad\BM.
\end{equation}
Sie ist genau dann $1$, wenn die Mode keine Wicklung in der Zeitrichtung hat ($k_4=0$), und für jede
massive Mode bei jedem Impuls kleiner als $1$; die Energie zum Erreichen von $v=1$ divergiert \BM.
Eine massive Mode verteilt ihre Wellenzahl auf Raum und $x_4$, und nur der Raumanteil ist Bewegung.
In der Massenlesart heißt das: Was eine Uhr trägt, ist langsamer als Licht; die Grenze erreicht nur,
was keine Eigenzeit hat (Dok.~391). Die kinetische Form von $\tilde T\cdot m=1$ ist de Broglies
Harmonie der Phasen: Die innere Frequenz sinkt um $\gamma$, die Phasenfrequenz steigt um $\gamma$, das
Produkt bleibt (Dok.~312) \BM.

Die Grenze ist in allen Darstellungen dieselbe. Eine konforme Skalierung $\Omega^2\eta_{\mu\nu}$ lässt
die Nullvektoren fest, und das lokal gemessene Verhältnis von Eigenlänge zu Eigenzeit eines
Lichtsignals hängt nicht von $\Omega$ ab \BM; die Massenvariation kann $c$ lokal nie ändern. Was sich
im Feld ändert, ist die Koordinatengeschwindigkeit: Sie ist bei reiner Massenvariation $1$ und im
isotropen Feld mit $\gamma=1$ gleich $1-2U=1/n$ mit $n=1+2U$, am Sonnenrand um $4{,}2\times10^{-6}$
kleiner \KM. Diese Verlangsamung gehört zur räumlichen Hälfte, der Wegverlängerung, nicht zu einem
veränderlichen $c$ (R132).

Dasselbe gilt für die elektrische und die magnetische Feldkonstante. Ein statisches Feld wirkt auf
Maxwell-Felder wie ein Medium mit $\varepsilon=\mu=\sqrt{B/A}$ für $g_{00}=-A$, $g_{ij}=B\,\delta_{ij}$
(Plebanski~\cite{Plebanski1960}) \QM. Bei $\gamma=1$ steigen beide gleich um $1+2U$, die
Koordinatengeschwindigkeit $1/\sqrt{\varepsilon\mu}$ sinkt um $2U$, und die Wellenimpedanz
$\sqrt{\mu/\varepsilon}=1$ bleibt \BM. Deshalb wird Licht im Schwerefeld weder reflektiert noch
doppelbrechend abgelenkt, und die Ablenkung hängt nicht von Farbe und Polarisation ab. Reine
Massenvariation lässt $\varepsilon=\mu=1$, weil die Maxwell-Gleichungen in vier Dimensionen konform
invariant sind \BM; auch hier trägt die räumliche Hälfte den ganzen Effekt, und zwar zu gleichen
Teilen in $\varepsilon$ und $\mu$. Lokal gemessen bleiben $\varepsilon_0$, $\mu_0$ und
$\varepsilon_0\mu_0c^2=1$ überall gleich; in Heaviside-Lorentz-Einheiten der Grundform sind beide eins.
Die Mitführung aus Abschnitt~6 erscheint in dieser Medium-Form als magnetoelektrische Kopplung
$\vec w=-g_{0i}/g_{00}$ zwischen elektrischem und magnetischem Feld (Plebanski) \QM.

% ============================================================
\section{Ergebnis}
% ============================================================

\begin{center}
\begin{tabular}{@{}p{4.2cm}p{5.4cm}p{3.6cm}@{}}
\toprule
Bereich & Massenlesart & Status \\
\midrule
Fall, Rotverschiebung, Uhren & trägt allein, gemeinsamer Kern & \KM \\
Universum im Großen & exakt (Weyl $=0$), besteht drei Diskriminatoren, AP wie FRW & \BM/\KM \\
Entfernungsgesetz $1+z=e^{\xi x}$ & AP-Test bei $z=2{,}33$ verfehlt ($5{,}3\,\sigma$) & \XM \\
Licht, Merkur, geodätisch & braucht die räumliche Hälfte, $\gamma=\beta=1$ & \KM \\
rotierende Massen & braucht ein mitgeführtes Vektorfeld & \BM, Herkunft \SM \\
Gravitationswellen & braucht einen Tensoranteil & \SM \\
Lichtgeschwindigkeit, $\varepsilon_0$, $\mu_0$ & lokal unverändert, Lichtkegel konform invariant & \BM \\
\bottomrule
\end{tabular}
\end{center}

Die drei Darstellungen aus Dok.~389 sind gleichwertig für den konform flachen Anteil und für das
statische, nicht rotierende Schwachfeld mit $\gamma=\beta=1$. Die Massenlesart ist die vollständige
Darstellung des konform flachen Anteils: Sie trägt die zeitliche Hälfte überall und das Universum im
Großen exakt. Was sich als Raumkrümmung im engeren Sinn zeigt, ist der Weyl-Anteil um einzelne
Massen; ihn muss die flache Darstellung zusätzlich liefern, als räumliche Hälfte mit dem
Maßstabsvorzeichen der Eigenlängen, als Mitführungsfeld bei Rotation und als Tensoranteil bei
Wellen. Mischformen sind zulässig, solange die räumliche Hälfte insgesamt den Anteil eins trägt.

% ============================================================
\section{Was offen bleibt}
% ============================================================

\begin{itemize}[leftmargin=*,itemsep=1pt]
\item Die Herkunft der räumlichen Hälfte aus der Struktur des Trägers. Dok.~308 weist die
K3-Flächeninvarianz als einzigen von vier geprüften Kandidaten aus, der $\gamma=1$ ergibt, und bestätigt
den Faktor zwei der Lichtablenkung numerisch; den Vorwärtsbeweis führt Dok.~308 selbst als offen,
ebenso Dok.~389 \SM.
\item Mitführung und Tensoranteil der Gravitationswellen in der flachen Darstellung. In der
Krümmungsdarstellung sind beide enthalten, weil die FFGFT die Feldgleichungen übernimmt, deren Form
nach Lovelock erzwungen ist (Dok.~312). Offen ist, ob die Struktur des Trägers das Mitführungsfeld
$\vec u$ und einen Tensoranteil selbst liefert. Der Korpus enthält dafür nur skalare Ansätze
(Dok.~131, 189); der Versuch in Dok.~143, den Tensoranteil aus dem Skalar zu gewinnen, ist nach
Dok.~144 zirkulär \SM.
\item Ein Entfernungsgesetz der FFGFT, das den Alcock--Paczyński-Test besteht. Die Reziprozität
$d_A=d_L/(1+z)^2$ ist nicht mehr offen: Sie gilt in der Lesart M exakt und ist mit $k=4$ dieselbe
Aussage. Die reine Exponentialform $1+z=e^{\xi x}$ (Dok.~280, 308) scheitert, als Koordinatenabstand
der Lesart M gelesen, am Lyman-$\alpha$-Wert von DESI (Abschnitt~5). Offen ist, welche Beziehung
$z(d)$ die Struktur des Trägers liefert und ob eine andere Lesung der Wegstrecke $x$ das Ergebnis
ändert; Dok.~268 führt $D_A$ bisher als externe Größe \SM.
\end{itemize}

\vspace{1em}
\noindent\textit{Prüfskripte:} \texttt{2/python/Dok392\_Skripte/} --- \texttt{pruef\_392a\_schwachfeld.py}
(14/14), \texttt{pruef\_392b\_mitfuehrung.py} (11/11), \texttt{pruef\_392c\_gravitationswellen.py} (8/8),
\texttt{pruef\_392d\_kosmos.py} (9/9), \texttt{pruef\_392e\_massenlesart.py} (11/11), \texttt{pruef\_392f\_lichtgrenze.py} (16/16), \texttt{pruef\_392g\_abstaende.py} (8/8),
\texttt{pruef\_392h\_desi\_ap.py} (4/4); zusammen \texttt{run\_all\_392.sh}, 81/81.

\section*{Werkzeug und Methode}
\addcontentsline{toc}{section}{Werkzeug und Methode}

Claude (Anthropic) wurde als Werkzeug eingesetzt für: Durchsuchen des Korpus, Schreiben und Ausführen
der Prüfskripte, Ausarbeiten der \LaTeX-Quellen nach den Vorgaben des Autors. Jede Zahl im Text ist
durch eine Prüfung in den genannten Skripten abgesichert; Angaben, die nur Literaturwerte wiedergeben,
sind als \QM\ markiert und werden dort nicht als Prüfung gezählt.

\begin{thebibliography}{99}
\bibitem{Bertotti2003} B.~Bertotti, L.~Iess, P.~Tortora, \textit{A test of general relativity using radio links with the Cassini spacecraft}, Nature \textbf{425} (2003) 374--376.
\bibitem{Everitt2011} C.~W.~F.~Everitt u.\,a., \textit{Gravity Probe B: Final results of a space experiment to test general relativity}, Phys.\ Rev.\ Lett.\ \textbf{106} (2011) 221101.
\bibitem{Williams2012} J.~G.~Williams, S.~G.~Turyshev, D.~H.~Boggs, \textit{Lunar laser ranging tests of the equivalence principle}, Class.\ Quantum Grav.\ \textbf{29} (2012) 184004.
\bibitem{Ciufolini2004} I.~Ciufolini, E.~C.~Pavlis, \textit{A confirmation of the general relativistic prediction of the Lense--Thirring effect}, Nature \textbf{431} (2004) 958--960.
\bibitem{Abbott2017a} B.~P.~Abbott u.\,a., \textit{Gravitational waves and gamma-rays from a binary neutron star merger: GW170817 and GRB 170817A}, Astrophys.\ J.\ Lett.\ \textbf{848} (2017) L13.
\bibitem{Abbott2017b} B.~P.~Abbott u.\,a., \textit{GW170814: A three-detector observation of gravitational waves from a binary black hole coalescence}, Phys.\ Rev.\ Lett.\ \textbf{119} (2017) 141101.
\bibitem{Weisberg2016} J.~M.~Weisberg, Y.~Huang, \textit{Relativistic measurements from timing the binary pulsar PSR B1913+16}, Astrophys.\ J.\ \textbf{829} (2016) 55.
\bibitem{Pound1960} R.~V.~Pound, G.~A.~Rebka, \textit{Apparent weight of photons}, Phys.\ Rev.\ Lett.\ \textbf{4} (1960) 337--341.
\bibitem{Chou2010} C.~W.~Chou u.\,a., \textit{Optical clocks and relativity}, Science \textbf{329} (2010) 1630--1633.
\bibitem{Vessot1980} R.~F.~C.~Vessot u.\,a., \textit{Test of relativistic gravitation with a space-borne hydrogen maser}, Phys.\ Rev.\ Lett.\ \textbf{45} (1980) 2081--2084.
\bibitem{Delva2018} P.~Delva u.\,a., \textit{Gravitational redshift test using eccentric Galileo satellites}, Phys.\ Rev.\ Lett.\ \textbf{121} (2018) 231101.
\bibitem{DESI2025} M.~Abdul Karim u.\,a.\ (DESI Collaboration), \textit{DESI DR2 Results II: Measurements of baryon acoustic oscillations and cosmological constraints}, arXiv:2503.14738 (2025), DOI 10.1103/tr6y-kpc6; Tabelle~IV.
\bibitem{Wetterich2013} C.~Wetterich, \textit{A Universe without expansion}, Phys.\ Dark Univ.\ \textbf{2} (2013) 184--187.
\bibitem{White2024} R.~M.~T.~White u.\,a., \textit{The Dark Energy Survey supernova program: slow supernovae show cosmological time dilation out to $z\sim1$}, Mon.\ Not.\ R.\ Astron.\ Soc.\ \textbf{533} (2024) 3365--3378.
\bibitem{Noterdaeme2011} P.~Noterdaeme u.\,a., \textit{The evolution of the cosmic microwave background temperature}, Astron.\ Astrophys.\ \textbf{526} (2011) L7.
\bibitem{Lubin2001} L.~M.~Lubin, A.~Sandage, \textit{The Tolman surface brightness test for the reality of the expansion. IV.}, Astron.\ J.\ \textbf{122} (2001) 1084--1103.
\bibitem{Plebanski1960} J.~Plebanski, \textit{Electromagnetic waves in gravitational fields}, Phys.\ Rev.\ \textbf{118} (1960) 1396--1408.
\end{thebibliography}
